% from ega3-6-fr.tex
\oldpage[III]{145}
En effet, vu l'exactitude du foncteur $M\otimes_A A'$ en $M$,
$L_{\bullet\bullet}\otimes_A A'$ et $M_{\bullet\bullet}\otimes_A A'$ sont
alors des résolutions projectives de Cartan-Eilenberg de
$P_\bullet\otimes_A A'$ et $Q_\bullet\otimes_A A'$ respectivement, d'où la
conclusion.
\end{proof}

\begin{env}[6.3.10]
Soit $\rho:A\to A'$ un homomorphisme d'anneaux; pour tout complexe
$P'_\bullet$ de $A'$-modules, $P'_{\bullet[\rho]}$ est un complexe de
$A$-modules; en outre, l'application identique
$P'_{\bullet[\rho]}\to P'_\bullet$ peut être considérée comme composée des
applications canoniques
\[
  P'_{\bullet[\rho]}\longrightarrow
  P'_{\bullet[\rho]}\otimes_A A'
  \xrightarrow{\mu}P'_\bullet,
\]
où $\mu$ est le $A'$-homomorphisme $\mu(x\otimes a')=a'x$. Si $Q'_\bullet$
est un second complexe de $A'$-modules, on a donc des homomorphismes
canoniques factoriels de degré 0
\[
  \mathbf{Tor}_\bullet^A(P'_{\bullet[\rho]},Q'_{\bullet[\rho]})
  \longrightarrow
  \mathbf{Tor}_\bullet^{A'}
  (P'_{\bullet[\rho]}\otimes_A A',Q'_{\bullet[\rho]}\otimes_A A')
  \longrightarrow
  \mathbf{Tor}_\bullet^{A'}(P'_\bullet,Q'_\bullet),
  \tag{6.3.10.1}\label{III.6.3.10.1.fr}
\]
où la première flèche est le $A$-homomorphisme défini dans (6.3.8) et la
seconde se déduit des $A'$-homomorphismes
$P'_{\bullet[\rho]}\otimes_A A'\to P'_\bullet$ et
$Q'_{\bullet[\rho]}\otimes_A A'\to Q'_\bullet$ par fonctorialité. On a des
homomorphismes analogues pour les suites spectrales de (6.3.2), et des
propriétés évidentes de transitivité, que nous laissons au lecteur le soin
d'énoncer.
\end{env}

\begin{proposition}[6.3.11]
Soit $\rho:A\to A'$ un homomorphisme d'anneaux faisant de $A'$ un
$A$-module plat. Pour tout complexe $P'_\bullet$ de $A'$-modules et tout
complexe $Q_\bullet$ de $A$-modules limités inférieurement, on a un
isomorphisme canonique factoriel
\[
  \mathbf{Tor}_\bullet^A(P'_{\bullet[\rho]},Q_\bullet)
  \xrightarrow{\sim}
  \mathbf{Tor}_\bullet^{A'}(P'_\bullet,Q_\bullet\otimes_A A').
  \tag{6.3.11.1}\label{III.6.3.11.1.fr}
\]
\end{proposition}

\begin{proof}
En effet, si $M_{\bullet\bullet}$ est une résolution projective de
Cartan-Eilenberg de $Q_\bullet$, $M_{\bullet\bullet}\otimes_A A'$ est une
résolution projective de Cartan-Eilenberg de $Q_\bullet\otimes_A A'$, et
l'on a, à un isomorphisme canonique près,
\[
  P'_{\bullet[\rho]}\otimes_A M_{\bullet\bullet}
  =P'_\bullet\otimes_{A'}(M_{\bullet\bullet}\otimes_A A');
\]
la conclusion résulte de (6.3.4).
\end{proof}

\begin{namedtheorem}[6.3.12]{Remarque}
Soit $(A^\lambda)$ un système inductif filtrant d'anneaux, et soient
$(P_\bullet^\lambda)$, $(Q_\bullet^\lambda)$ deux systèmes inductifs de
complexes de $(A^\lambda)$-modules; on a alors un isomorphisme canonique
généralisant (6.3.6.1)
\[
  \varinjlim_\lambda
  \mathbf{Tor}_\bullet^{A^\lambda}
  (P_\bullet^\lambda,Q_\bullet^\lambda)
  \xrightarrow{\sim}
  \mathbf{Tor}_\bullet^A(P_\bullet,Q_\bullet),
  \tag{6.3.12.1}\label{III.6.3.12.1.fr}
\]
où $A=\varinjlim A^\lambda$, $P_\bullet=\varinjlim P_\bullet^\lambda$,
$Q_\bullet=\varinjlim Q_\bullet^\lambda$. Une fois définis les
homomorphismes
\[
  \mathbf{Tor}_\bullet^{A^\lambda}
  (P_\bullet^\lambda,Q_\bullet^\lambda)
  \longrightarrow
  \mathbf{Tor}_\bullet^{A^\mu}
  (P_\bullet^\mu,Q_\bullet^\mu)
\]
pour $\lambda\leqslant\mu$, à l'aide de (6.3.10), la démonstration est celle
de (6.3.6).
\end{namedtheorem}

\begin{proposition}[6.3.13]
Soient $S$ une partie multiplicative de $A$, $P_\bullet$ et $Q_\bullet$ deux
complexes de $A$-modules, dans lesquels les homothéties définies par les
éléments de $S$ soient bijectives, de sorte que, si $A'=S^{-1}A$,
$P_\bullet$ et $Q_\bullet$ sont formés de $A'$-modules. Alors on a un
isomorphisme canonique
\[
  \mathbf{Tor}_\bullet^A(P_\bullet,Q_\bullet)
  \xrightarrow{\sim}
  \mathbf{Tor}_\bullet^{A'}(P_\bullet,Q_\bullet).
\]
\end{proposition}

\begin{proof}
En effet, l'hypothèse entraîne que les homomorphismes canoniques
$P_\bullet\to P_\bullet\otimes_A A'$,
$Q_\bullet\to Q_\bullet\otimes_A A'$ sont bijectifs. D'autre part, la
fonctorialité de l'hypertor montre que tout $s\in S$ définit une homothétie
bijective dans $\mathbf{Tor}_\bullet^A(P_\bullet,Q_\bullet)$, et par suite
\[
  \mathbf{Tor}_\bullet^A(P_\bullet,Q_\bullet)
  \longrightarrow
  \mathbf{Tor}_\bullet^A(P_\bullet,Q_\bullet)\otimes_A A'
\]

